% titel: Blatt 4, A1: Photonen (Lampe) und Aussagen zur Elektronenstruktur
% Teil II: der durchgerechnete Weg.

\subsection*{A1.1 --- Zahl der Photonen}
Eine Lampe emittiert bei $\lambda = \qty{450}{nm}$ mit einer Leistung von
$P = \qty{5.0}{J/s}$. Gesucht: Photonen in $t = \qty{12.5}{s}$.

\subsection*{Rechnung}
Energie eines Photons:
\begin{equation*}
  E_{\text{photon}} = \frac{h\,c}{\lambda}
  = \frac{6.626\times10^{-34}\times 2.998\times10^{8}}
       {4.50\times10^{-7}}
  = 4.42\times10^{-19}\ \mathrm{J}.
\end{equation*}
Energie der Lampe in 12.5 s:
\begin{equation*}
  E = P\,t = 5.0 \times 12.5 = \qty{62.5}{J}.
\end{equation*}
Anzahl der Photonen:
\begin{equation*}
  N = \frac{E}{E_{\text{photon}}}
  = \frac{62.5}{4.42\times10^{-19}}
  = 1.41\times10^{20}
  \approx \num{1.4e20}.
\end{equation*}
Die richtige Antwort ist also der Wert $\approx \num{1.4e20}$
(Option ``1.4$\,\times\,10^{20}$'').

\subsection*{A1.2 --- Aussagen wahr/falsch}
\begin{itemize}
  \item ``Rückfall von angeregt in den Grundzustand gibt Energie frei,
        das ist Emission'': \textbf{wahr} ($+$).
  \item ``Heisenberg: Ort und Impuls nicht gleichzeitig genau bestimmbar'':
        \textbf{wahr} ($+$).
  \item ``Elektronen halten sich überall im Atom gleich wahrscheinlich
        auf'': \textbf{falsch} ($-$) --- die Aufenthaltswahrscheinlichkeit
        ist durch die Orbitale räumlich strukturiert.
  \item ``Modell ``Elektron im Kasten'' erklärt die Energiequantisierung'':
        \textbf{wahr} ($+$).
\end{itemize}

\begin{ergebnis}
A1.1: $N \approx \num{1.4e20}$ Photonen. A1.2: Antwortfolge
$+$, $+$, $-$, $+$.
\end{ergebnis}

\begin{kontrolle}
Energiebilanz: $62.5\ \mathrm{J}$ geteilt durch
$\approx 4.4\times10^{-19}\ \mathrm{J/Photon}$ liefert größenordnungsmäßig
$10^{20}$ --- konsistent mit sichtbarem Licht bei einigen Joule. Die
richtige Option stimmt mit der Rechnung überein, die drei anderen (13, 15,
16) liegen falsch.
\end{kontrolle}

\begin{merksatz}
$E_{\text{photon}} = hc/\lambda$, $E = P\,t$, $N = E/E_{\text{photon}}$;
Konzeptfragen testen vor allem: Emission, Heisenberg und die Tatsache, dass
die Elektronenaufenthaltswahrscheinlichkeit nicht überall gleich ist.
\end{merksatz}
